\chapterspace{-2}
\chapter{Ejemplos y complementos}
\label{XI}
\marginpar{285}

\section{Espacios proyectivos, variedades unirracionales}
\label{XI.1}

\begin{proposition}
\label{XI.1.1}
Sean $k$ un cuerpo algebraicamente cerrado y $X=\PP_k^r$ el espacio
proyectivo de dimensión $r$ sobre~$k$. Entonces $X$ es \emph{simplemente
conexo,} es decir, $\pi_1(X)=0$.
\end{proposition}

Para $r=0$, esto es trivial. Si $r=1$, hay que demostrar que, si $X'$ es
un recubrimiento étale conexo no vacío de~$X=\PP_k^1$, entonces
$X'\isomto X$. La fórmula del género nos da aquí, si~$g$ y $g'$ son
los géneros de~$X$ y~$X'$:
$$
1-g' = d(1-g)\quoi,
$$
donde $d$ es el grado de~$X'$ sobre~$X$. Como $g=0$, se tendrá entonces
$1-g'=d$, lo cual exige $d=1$, puesto que $g'\ge 0$; esto demuestra
$X'\isomto X$. Cuando $r\ge 2$, se procede por inducción
sobre~$r$, suponiendo que $\PP^{r'}$ es simplemente conexo para
$r'<r$. Aplicando esto a un hiperplano de~$\PP^r$ y utilizando
(\Ref{X.2.10}), resulta que $\PP^r$ es simplemente conexo. Otra
demostración: en virtud de (\Ref{X.1.7}), se tiene
$\pi_1(\PP^1\times\cdots\times \PP^1) = \pi_1(\PP^1)\times\cdots
\times \pi_1(\PP^1)$; por tanto, $(\PP^1)^r$ es simplemente conexo,
pues $\PP^1$ lo es, y entonces $\PP^r$ es simplemente conexo en virtud
de la invariancia birracional del grupo fundamental (\Ref{X.3.4}). Esta
demostración prueba, más generalmente:

\begin{corollaire}
\label{XI.1.2}
Sea $X$ un esquema propio y normal sobre un cuerpo algebraicamente
cerrado~$k$; si $X$ es una variedad racional, es decir, si es íntegro y
su cuerpo de funciones es una extensión puramente trascendente de~$k$,
entonces $X$ es simplemente conexo.
\end{corollaire}

Este resultado se aplica, en particular, a las variedades grassmannianas y,
más
\marginpar{286}%
generalmente, a las variedades~$G/H$, donde $G$ es un grupo lineal conexo
sobre~$k$ y $H$ un subgrupo algebraico que contiene un subgrupo de Borel
\ifthenelse{\boolean{orig}}
{de~$B$.}
{de~$G$.}

Recordemos que se llama variedad \emph{unirracional sobre~$k$}
\index{variedad unirracional|hyperpage}%
a un esquema propio e íntegro sobre~$k$ cuyo cuerpo de funciones $K$
está contenido en una extensión puramente trascendente~$K'$ de~$k$,
finita sobre~$K$ (es decir, que tiene el mismo grado de trascendencia
sobre~$k$ que~$K$), o, equivalentemente, si existe una aplicación racional
dominante $f\colon\PP_k^r\to X$, con $r=\dim X$. Si $X$ es normal,
se deduce de las consideraciones que preceden a \Ref{X.3.4} que, para todo
recubrimiento étale conexo $X'$ de~$X$, de cuerpo~$L/K$, el álgebra
$L\otimes_K K'$ sobre~$K'$ no está ramificada sobre el modelo~$\PP^r$ y,
por tanto, está completamente descompuesta en virtud de~\Ref{X.1.1}; esto
muestra que $L$ es $K$-isomorfo a una subextensión de~$K'/K$. Teniendo en
cuenta \Ref{V.8.2}, queda demostrado:

\begin{corollaire}
\label{XI.1.3}
El grupo fundamental de una variedad unirracional normal sobre un cuerpo
algebraicamente cerrado es finito.
\end{corollaire}

(N.B. Obsérvese que, en la definición de «unirracional», no era necesario
suponer que $K'/K$ fuera finita).

\ifthenelse{\boolean{orig}}
{\begin{remarques}
\label{XI.1.4}
Naturalmente, los resultados de este número son bien conocidos. Por otra
parte,
\ifthenelse{\boolean{orig}}{J.P\ptbl}{J-P\ptbl}%
Serre demostró~\cite{XI.10} que, cuando $X$ es una variedad proyectiva
unirracional lisa sobre un cuerpo algebraicamente cerrado de
\emph{característica cero,} $X$ es simplemente conexa. Su demostración es
trascendente, en el sentido de que utiliza el teorema de simetría de Hodge,
y se ignora si este resultado se extiende a la característica~$p>0$.
Parece, además, que no se conoce ningún ejemplo de variedad unirracional
lisa sobre~$k$ que no sea ya racional.
\end{remarques}}
{\begin{remarques}
\label{XI.1.4}
Naturalmente, los resultados de este número son bien conocidos. Por otra
parte,
\ifthenelse{\boolean{orig}}{J.P\ptbl}{J-P\ptbl}%
Serre demostró~\cite{XI.10} que, cuando $X$ es una variedad proyectiva
unirracional lisa sobre un cuerpo algebraicamente cerrado de
\emph{característica cero,} $X$ es simplemente conexa. Su demostración es
trascendente, en el sentido de que utiliza el teorema de simetría de Hodge.

\lcrochetbf Añadido en 2003 (MR) Limitémonos al caso de un cuerpo
algebraicamente cerrado $k$ de característica $p\geq 0$. En característica
$p>0$, la definición de $k$-variedad unirracional dada más arriba es la de
$k$-variedad débilmente unirracional, por oposición a la de $k$-variedad
fuertemente unirracional, en la que se supone además que $K'$ es una
extensión separable de~$K$.

En dimensión $2$, existen superficies proyectivas y lisas, débilmente
unirracionales, cuyo grupo fundamental (finito por \Ref{XI.1.3}) no es
trivial y que, por tanto, no son racionales (T\ptbl Shioda, On
unirationality of supersingular surfaces, Math.~Ann. \textbf{225} (1977),
p\ptbl155--159.). En cambio, una superficie fuertemente unirracional es
racional: esto resulta del criterio de racionalidad de Castelnuovo,
extendido a la característica $p>0$ por O\ptbl Zariski (\cf J-P\ptbl Serre,
Sém. Bourbaki \no146, 1957 y {\OE}uvres--Collected papers, vol\ptbl 1,
p\ptbl 467).

Sobre el cuerpo $\CC$ de los números complejos, se conocen ejemplos de
variedades proyectivas y lisas de \hbox{dimensión $\geq 3$} que son
unirracionales y no racionales (\cf P\ptbl Deligne, Variétés
unirationnelles non rationnelles [d'après M\ptbl Artin et D\ptbl Mumford],
Sém.\ Bourbaki \no402, 1971-72, Lecture Notes vol\ptbl 317). Este es, en
particular, el caso de una hipersuperficie cúbica lisa del espacio
proyectivo $\PP^4$. Algunos de estos ejemplos se extienden a característica
$>0$ y proporcionan variedades fuertemente unirracionales no racionales
(\cf J.P\ptbl Murre, Reduction of the proof of the non-rationality of a non
singular cubic threefold to a result of Mumford, Compositio Math.
\textbf{27} (1973), p\ptbl63--82).

Sea $V$ una $k$-variedad proyectiva, normal e íntegra. Se dice que $V$ es
racionalmente conexa si existen un $k$-esquema íntegro de tipo finito $T$ y
un $k$-morfismo $F : T\times\PP^1\to V$ tales que el morfismo
\begin{equation}\tag*{(1)}
\begin{split}
T \times\PP^1\times\PP^1&\to V\times V\\
(t, u, u')&\mto(F(t,u), F(t,u'))
\end{split}
\end{equation}
sea dominante. En particular, por dos puntos racionales suficientemente
generales de~$V$ pasa una curva racional. La terminología se justifica por
el hecho de que, si $V$ es racionalmente conexa, se pueden unir dos puntos
racionales mediante una cadena finita de curvas racionales. Si $k$ tiene
característica $p>0$, se refuerza la definición anterior exigiendo que la
aplicación (1) sea genéricamente lisa. Se obtiene así la noción de variedad
racionalmente conexa separable. Por tanto, las variedades unirracionales son
racionalmente conexas y, en característica $p>0$, las variedades
fuertemente unirracionales son racionalmente conexas separables. J\ptbl
Kollár demostró que las variedades racionalmente conexas separables tienen
grupo fundamental algebraico trivial (\cf O\ptbl Debarre, Variétés
rationnellement connexes [d'après T\ptbl Graber, J\ptbl Harris, J\ptbl
Starr et A.J\ptbl de~Jong], Sém. Bourbaki \no906, 2001-2002).\rcrochetbf
\end{remarques}}
